\section{Reacciones orgánicas hmolíticas y heterolíticas}

\subsection{Reacciones homolítica}

La ruptura del enlace covalente ocurre de manera simétrica, es decir, cada átomo involucrado en el enlace se lleva consigo su electrón

Características:

\begin{itemize}
    \item Ruptura del enlace simétrica
    \item Formación de radicales libres
    \item Generalmente iniciada por energía
\end{itemize}

\subsection{Reacciones heterolítica}

Implica un rompimiento de un enlace covalente. Uno de los átomos se lleva consigo el par electrónico, generando iones.

Características:
\begin{itemize}
    \item Ruptura asimétrica de los enlaces
    \item Formación de iones
    \item Influenciada por la polaridad y naturaleza del solvente
\end{itemize}

\section{Cloración del metano}

La cloración del metano es una reacción química entre el cloro y el metano la cual produce una mezcla de productos clorados. Se caracteriza además por ser una reacción en cadena que da lugar a múltiples productos a partir de la misma aplicación de energía

Las características principales de esta reacción pueden ser englobadas en los siguientes 3 puntos:

\begin{itemize}
    \item La cloración no ocurre a temperatura ambiente en ausencia de luz. Esto porque la reacción necesita de una fuente de energía para ocurrir.
    \item La longitud de onda de la luz que resulta más eficaz es la de color azul, ya que es fuertemente absorbida por el cloro en estado gaseoso. Esto es a causa de que su energía (242 kJ/mol) es la más adecuada para disociar las moléculas de cloro (C2) y que así quede un cloro libre para interactuar con el metano.
    \item La reacción iniciada por la luz tiene un rendimiento cuántico elevado. Sucede así porque la cloración del metano tiene un mecanismo de reacción en cadena, pues cada cloración deja libre un intermediario reactivo que puede generar otra reacción con otra molécula de metano, por lo que una sola aplicación de energía, como un fotón, puede dar pie a más de una reacción de cloración del metano.
\end{itemize}

\subsection{Pasos de propagación}

Esta reacción tiene distintos pasos, los cuales son:

\begin{enumerate}
    \item Fase de iniciación, donde se generan los intermediarios reactivos. El comienzo de la reacción, dado en la fase de iniciación, se da cuando a las moléculas de cloro se les aplica energía, ya sea calor o luz. En el caso de la aplicación de luz, es necesario que esta sea luz azul (400-495 nm), puesto que la cantidad de energía que proporciona es la adecuada para que las moléculas de cloro (Cl2) se disocien entre sí a través de una ruptura homolítica y creen un par de intermediarios reactivos, llamados también radicales libres por contar con electrones desapareados. Posterior a esto sigue el paso de propagación, el cual consta de 2 pasos: la generación de un radical metilo y la asociación de este con un átomo de cloro. En el primero, un radical libre de cloro choca con una molécula de metano y, por consecuencia, retira a uno de los átomos de hidrógeno enlazados al metano. Al ser esta una ruptura homolítica, el hidrógeno saliente conserva su electrón de valencia y se enlaza a través de este con el radical claro para convertirse en cloruro de hidrógeno (ClH), el cual ya no reacciona con otras moléculas en este proceso.
    
    \item Fase de propagación, donde los intermediarios reactivos interactúan con las moléculas de metano para generar los productos de la reacción. El segundo paso ocurre cuando el radical metilo choca ahora con una molécula de cloro (Cl2). Al suceder esto, se rompe el enlace entre los átomos de cloro por medio de homólisis y uno de estos pasa a enlazarse con el radical metilo a través de un enlace covalente entre sus electrones impares. El cloro que no realizó el enlace continúa en el medio hasta que choca con otra molécula de metano y produce un nuevo radical metilo, dando inicio otra vez al segundo paso de propagación.
        
    \item Fase de terminación, donde la reacción no cuenta con más reactivos para continuar la reacción en cadena. Este proceso debe ser constante, pues cuando algún radical metilo se une a otro reactivo diferente al cloro (o viceversa) se rompe la cadena de propagación, por lo que el ritmo de la reacción comienza a decrecer hasta llegar a 0, osea, la reacción termina. A este rompimiento de la cadena de reacción es a lo que se conoce como fase de terminación. Algunos de los reactivos a los que se podría unir el radical metilo para romper tal cadena podrían ser:
        \begin{itemize}
            \item Otro radical metilo, formando etano. La propia pared del contenedor en que se realice la reacción o algún contaminante presente en la muestra.
            \item El cloro, por su parte, rompería la cadena de reacción si se enlazara con lo siguiente: Otro átomo de cloro, produciendo de nuevo una molécula de cloro (Cl2).
            \item Con la propia pared del contenedor o algún contaminante presente o algún contaminante presente en la muestra, al igual que el radical metilo. Las probabilidades de que uno de los radicales reaccione con lo antes mencionado y que por ende comience la fase de terminación aumenta con el tiempo a causa de la reducción en la cantidad de radicales presentes en el medio total.
        \end{itemize}
\end{enumerate}
